\documentclass[11pt]{article}
\usepackage[T1]{fontenc}
\usepackage[utf8]{inputenc}
\usepackage{lmodern,microtype}
\usepackage[margin=1.08in]{geometry}
\usepackage{amsmath,amssymb,amsthm,mathtools}
\usepackage{xcolor,hyperref}
\hypersetup{colorlinks=true,linkcolor=blue!55!black,citecolor=blue!55!black,urlcolor=blue!60!black}
\setlength{\parindent}{0pt}
\setlength{\parskip}{0.48em}
\setlength{\emergencystretch}{2em}
\newtheorem{theorem}{Theorem}
\newtheorem{lemma}[theorem]{Lemma}
\theoremstyle{remark}
\newtheorem{remark}[theorem]{Remark}

\title{A Nonconvex Hyperbolic Obstacle with Positive Exterior Robin Threshold}
\author{[Author Name]}
\date{}

\begin{document}
\maketitle

\begin{abstract}
\textbf{Relation to the problem list.}
This manuscript addresses Question~41 (L41), ``Positive exterior-Robin
threshold in the hyperbolic plane.''  It gives a full construction and proof
in the Lipschitz category stated in the source.  The strict spectral
inequality also survives rounding the four corners, so the example can be
made smooth.

We construct a bounded, nonconvex, simply connected Lipschitz obstacle in
the hyperbolic plane whose exterior Robin Laplacian has a discrete eigenvalue
below the essential-spectrum threshold $1/4$ for a positive Robin parameter.
The obstacle is an annular sector obtained by cutting a thin radial channel
through a geodesic annulus.  A compactly supported trial function, constant
in the inner chamber and decaying inside the channel, gives an explicit
strict Rayleigh bound.  In particular, the critical exterior Robin
parameter of the constructed obstacle is positive.
\end{abstract}

\section{Introduction}

Let $\Omega\subset\mathbb H^2$ be a bounded Lipschitz obstacle with
connected exterior
\[
 D=\mathbb H^2\setminus\overline\Omega.
\]
For $\alpha\in\mathbb R$, consider the closed quadratic form
\begin{equation}\label{eq:form}
 Q_\alpha^D[u]
 =\int_D|\nabla u|^2\,d\mu
 +\alpha\int_{\partial D}|u|^2\,d\sigma,
 \qquad u\in W^{1,2}(D).
\end{equation}
It defines the exterior Robin Laplacian in the sign convention of
Celentano, Krej\v{c}i\v{r}\'{i}k, and Lotoreichik \cite{CKL}.  Its lowest
spectral point is
\begin{equation}\label{eq:rayleigh}
 \lambda_1^\alpha(D)
 =\inf_{0\ne u\in W^{1,2}(D)}
 \frac{Q_\alpha^D[u]}{\|u\|_{L^2(D)}^2}.
\end{equation}
For bounded Lipschitz obstacles with connected exterior,
\begin{equation}\label{eq:ess}
 \sigma_{\mathrm{ess}}(-\Delta_\alpha^D)=[1/4,\infty)
\end{equation}
for every $\alpha\in\mathbb R$ \cite[Proposition~3.4]{CKL}.  Thus a trial
function with Rayleigh quotient below $1/4$ produces a discrete eigenvalue.

Define the critical threshold
\[
 \alpha_*(D)=\sup\{\alpha\in\mathbb R:\lambda_1^\alpha(D)<1/4\}.
\]
Our main result gives an explicit family for which this threshold is positive.

\begin{theorem}\label{thm:main}
Fix $0<r<R$ and set
\[
 A_r=2\pi(\cosh r-1),
 \qquad
 P_r=2\pi\sinh r,
 \qquad
 \alpha_r=\frac18\tanh\frac r2>0.
\]
For all sufficiently small $\eta>0$, the domain
\begin{equation}\label{eq:obstacle}
 \Omega_\eta
 =\{(\rho,\theta):r<\rho<R,\ \eta<\theta<2\pi-\eta\}
 \subset\mathbb H^2
\end{equation}
is a bounded, nonconvex, simply connected Lipschitz obstacle with connected
exterior $D_\eta$, and
\[
 \lambda_1^{\alpha_r}(D_\eta)<\frac14.
\]
Consequently, $-\Delta_{\alpha_r}^{D_\eta}$ has a discrete eigenvalue below
$1/4$ and $\alpha_*(D_\eta)>0$.
\end{theorem}

\section{Geometry of the obstacle}

Use geodesic polar coordinates about a point $o\in\mathbb H^2$:
\[
 g=d\rho^2+\sinh^2\rho\,d\theta^2,
 \qquad
 d\mu=\sinh\rho\,d\rho\,d\theta.
\]
The coordinate rectangle in \eqref{eq:obstacle} has four transverse
piecewise smooth corners.  Hence $\Omega_\eta$ is bounded, Lipschitz, and
diffeomorphic to a rectangle, so it is simply connected.  Its exterior is
the union of the inner disk, the removed radial wedge, and the region outside
the radius-$R$ disk.  The wedge joins these pieces, and therefore $D_\eta$
is connected.

If $0<\eta<\pi/2$, choose two points of the same radius
$s\in(r,R)$ and polar angles $\pi/2$ and $3\pi/2$.  Their unique minimizing
geodesic is the diameter through $o$, which crosses the disk $B_r(o)$.
Since this disk is disjoint from $\Omega_\eta$, the obstacle is not
geodesically convex.

\section{The trial function}

Let $\delta=\sqrt\eta<R-r$ and put $\chi(t)=(1-t)_+$.  Write $W_\eta$ for
the angular wedge of total aperture $2\eta$ centered on the ray $\theta=0$.
On $D_\eta$ define
\begin{equation}\label{eq:trial}
 u_\eta(\rho,\theta)=
 \begin{cases}
  1,&0\le\rho<r,\\[1mm]
  \chi\!\left(\dfrac{\rho-r}{\delta}\right),
  &r\le\rho<r+\delta,\quad\theta\in W_\eta,\\[3mm]
  0,&\text{otherwise}.
 \end{cases}
\end{equation}
The values match across the mouth $\rho=r$ and across
$\rho=r+\delta$.  The two angular sides of the decay region lie on
$\partial D_\eta$, so there is no interior jump there.  Standard Sobolev
gluing gives $u_\eta\in W^{1,2}(D_\eta)$.

\begin{lemma}\label{lem:estimates}
The trial function \eqref{eq:trial} satisfies
\begin{align}
 \int_{D_\eta}|\nabla u_\eta|^2\,d\mu
 &\le2\delta\sinh(r+\delta),\label{eq:energy}\\
 \int_{\partial D_\eta}|u_\eta|^2\,d\sigma
 &\le P_r+2\delta,\label{eq:trace}\\
 \|u_\eta\|_{L^2(D_\eta)}^2&\ge A_r.\label{eq:norm}
\end{align}
\end{lemma}

\begin{proof}
Only the radial decay region contributes to the gradient.  Since it has
angular aperture $2\eta$ and $\eta=\delta^2$,
\begin{align*}
 \int_{D_\eta}|\nabla u_\eta|^2\,d\mu
 &=\frac{2\eta}{\delta^2}
   \int_r^{r+\delta}\sinh\rho\,d\rho\\
 &=2\bigl(\cosh(r+\delta)-\cosh r\bigr)
 \le2\delta\sinh(r+\delta).
\end{align*}

The trace equals one on the inner circular arc, vanishes on the outer arc,
and equals $\chi((\rho-r)/\delta)$ on each of the two radial walls.  Hence
\begin{align*}
 \int_{\partial D_\eta}|u_\eta|^2\,d\sigma
 &=(2\pi-2\eta)\sinh r
 +2\int_r^{r+\delta}
 \left(1-\frac{\rho-r}{\delta}\right)^2d\rho\\
 &=(2\pi-2\eta)\sinh r+\frac{2\delta}{3}
 \le P_r+2\delta.
\end{align*}
Finally, $u_\eta=1$ throughout the inner disk $B_r(o)$, whose area is
$A_r$, proving \eqref{eq:norm}.
\end{proof}

\section{Proof of the spectral inequality}

\begin{proof}[Proof of Theorem~\ref{thm:main}]
By \eqref{eq:rayleigh} and Lemma~\ref{lem:estimates},
\begin{equation}\label{eq:quotient}
 \lambda_1^{\alpha_r}(D_\eta)
 \le
 \frac{\alpha_rP_r
2\delta\bigl(\sinh(r+\delta)+\alpha_r\bigr)}{A_r}.
\end{equation}
The hyperbolic disk identities give
\[
 \frac{P_r}{A_r}
 =\frac{\sinh r}{\cosh r-1}
 =\coth\frac r2,
 \qquad
 \frac{\alpha_rP_r}{A_r}=\frac18.
\]
The error term in \eqref{eq:quotient} tends to zero with
$\delta=\sqrt\eta$.  Thus the right-hand side is less than $1/4$ for all
sufficiently small $\eta$.

For an explicit range, put
\[
 C_r=\sinh\left(r+\frac14\right)+\alpha_r
\]
and choose
\[
 0<\eta<
 \min\left\{
 \frac\pi2,\frac1{16},(R-r)^2,
 \left(\frac{A_r}{16C_r}\right)^2
 \right\}.
\]
Then $\delta<1/4$, $\delta<R-r$, and
\[
 \frac{2\delta(\sinh(r+\delta)+\alpha_r)}{A_r}
 \le\frac{2\delta C_r}{A_r}<\frac18.
\]
Equation \eqref{eq:quotient} therefore yields
$\lambda_1^{\alpha_r}(D_\eta)<1/4$.  By \eqref{eq:ess}, the spectrum below
$1/4$ is discrete.  Since $\alpha_r>0$ belongs to the defining set for
$\alpha_*(D_\eta)$, one has
$\alpha_*(D_\eta)\ge\alpha_r>0$.
\end{proof}

\begin{remark}
The theorem is stated in the Lipschitz category used in \cite{CKL}.  Because
the Rayleigh inequality is strict, the four corners can also be rounded in
small disjoint neighborhoods while preserving the conclusion.
\end{remark}

\begin{thebibliography}{9}

\bibitem{CKL}
A.~Celentano, D.~Krej\v{c}i\v{r}\'{i}k, and V.~Lotoreichik,
\emph{Optimisation of the lowest Robin eigenvalue in exterior domains of the
hyperbolic plane},
\href{https://arxiv.org/abs/2601.10280}{arXiv:2601.10280}.

\bibitem{Chavel}
I.~Chavel,
\emph{Eigenvalues in Riemannian Geometry},
Academic Press, Orlando, 1984.

\end{thebibliography}
\end{document}
